% --- Back cover ---
% Large PERSONIX wordmark, no logo image. Cream paper inherited globally.
% Absolute (tikz overlay) positioning so the footer sticks to the bottom.
\clearpage
\thispagestyle{empty}
\begin{tikzpicture}[remember picture, overlay]
  % --- Wordmark ---
  \node[anchor=north, font=\sffamily\bfseries\fontsize{48}{52}\selectfont]
    at ([yshift=-26mm]current page.north) {PERSONIX};
  \node[anchor=north, font=\rmfamily\itshape\large]
    at ([yshift=-44mm]current page.north) {Bezkompromisa pārmaiņas};
  \node[anchor=north, font=\rmfamily\small,
        text width=\paperwidth-50mm, align=center]
    at ([yshift=-54mm]current page.north)
    {Necenzējams un neuzpērkams decentralizēts reputācijas tīkls};
  % --- Body ---
  \node[anchor=north, text width=\paperwidth-50mm, align=justify, font=\small]
    at ([yshift=-72mm]current page.north)
    {%
      Valsts darbojās, kad komunikācija bija lēna, lēmumi bija lokāli un
      autoritāte dzīvoja tajā pašā pilsētā, kuru pārvaldīja. Mūsdienu tīkli
      apgriež otrādi visus trīs pieņēmumus --- un uz tiem uzbūvētās
      institūcijas vairs neatbilst pasaulei, ko tās pārvalda. Šī grāmata
      apraksta rīku, kas atbilst: decentralizētu reputācijas tīklu, kurā
      identitāte pieder tev, patiesība ir publiski auditējama un kukuļošana
      ir reputācijas pašnāvība.

      \smallskip
      Ne manifests. Ne prognoze. Darbspējīgs plāns tam, kā uzbūvēt valsts
      pēcteci --- brīvprātīgi, pa vienai deklarācijai.
    };
  % --- Footer (anchored to the bottom of the page) ---
  \node[anchor=south, font=\sffamily\footnotesize,
        text width=\paperwidth-50mm, align=center]
    at ([yshift=12mm]current page.south)
    {%
      Pavel Kudrna\quad\textbullet\quad Prāga, 2026\par
      \href{https://personix.org}{personix.org}\par
      Licencēts saskaņā ar Creative Commons Attribution-ShareAlike 4.0 International
      (CC BY-SA 4.0).
    };
\end{tikzpicture}%
\mbox{}%
\clearpage
